\documentclass[10pt,a4paper]{amsart}
\usepackage[margin=19mm]{geometry}
\usepackage{amsmath,amssymb,mathtools}
\usepackage[scr=rsfs,cal=euler]{mathalfa}
\usepackage{xfrac}
\usepackage[unicode,hidelinks,bookmarks=false]{hyperref}
\newcommand{\ie}{i.e.\ }
\newcommand{\cf}{cf.\ }
\newcommand{\resp}{respectively\ }
\newcommand{\Subsection}{\subsection}
\providecommand{\nobreakdash}{\nobreak\hbox{-}}
\setlength{\emergencystretch}{2em}
\multlinegap0pt
\usepackage{amssymb}
\usepackage{enumerate}
\usepackage[shortlabels]{enumitem}
\usepackage{tikz}
\usepackage{float}
\newtheorem{lemma}{Lemma}
    \newcommand{\Br}{{\operatorname{Br}}}
\title{Towards Monoidal Categorifications of Twisted Products of Flag Varieties}
\author{Yingjin Bi}
\date{Source: arXiv:2602.11559v6 (CC BY 4.0)}
\begin{document}
\maketitle

\noindent\emph{Source and licence.} Towards Monoidal Categorifications of Twisted Products of Flag Varieties. Version arXiv:2602.11559v6 (CC BY 4.0), distributed by arXiv under the Creative Commons Attribution 4.0 International licence, http://creativecommons.org/licenses/by/4.0/, as recorded in the arXiv licence metadata for this exact version and shown on the abstract page https://arxiv.org/abs/2602.11559v6. Full licence notice: \texttt{licenses/arxiv-2602.11559-CC-BY-4.0.txt}.\par\smallskip

\noindent\textit{Editorial scope.} Counted unit: the article's lemma lem:beta3move (source0.tex lines 3654--3733). The article's complete original proof of it is delivered with it (source0.tex lines 3735--3925, one \texttt{proof} environment of 191 lines). No mathematics is written, repaired or re-derived: the statement and the proof are the article's own lines, in the article's own notation, and no retained line is substituted or re-flowed.  Retained uncounted as the source-local context the proof needs: lines 1308--1315.  Nothing was omitted from any retained span, so the package carries no omission record.  No retained line contains a citation, so the wrapper carries no bibliography and no bibliography entry.  No retained line contains a figure environment, an image-inclusion command or drawing code, so no image is delivered and the package declares no asset dependency.  Editorial formatting only, in addition: an AMS \texttt{amsart} preamble that restates the article's own declarations and macros used by the retained lines, namely the environments \texttt{lemma} on separate counters, together with the standard packages those lines need.  The retained environments print in this document as Lemma 1 in order of appearance, whereas the article prints the same items with the numbering of its own full document.  The complete native source of the article as distributed in the arXiv e-print for this exact version is bundled unchanged as \texttt{sources/194480/source0.tex}.
\begin{lemma}\label{lem:vk}
Let \(\beta=(i_1,\dots,i_r)\) be an expression of \(b\), and let
\(v\le \delta(b)\). Then
\[
v_k=v'_k
\]
for all \(0\le k\le r\).
\end{lemma}

\begin{lemma}\label{lem:beta3move}
Let \(v\le \delta(b)\), and let
\(
\beta=(i_1,\dots,i_r)
\)
be an expression of \(b\in \Br^+\). Let
\(
\beta_v=(p_1,\dots,p_m)
\)
be the sequence of positions of the leftmost reduced subexpression of \(v\) in
\(\beta\), where \(m=\ell(v)\). Then the following hold.

\begin{enumerate}
\item For \(i\in I\), one has
\[
s_iv<v
\quad\Longleftrightarrow\quad
v^{-1}(\alpha_i)\in R^-.
\]

\item Suppose that \(\beta'\) is obtained from \(\beta\) by a \(3\)-move
\[
(i_{k-1},i_k,i_{k+1})=(i,j,i)
\quad\longleftrightarrow\quad
(j,i,j),
\]
where \(c_{ij}=-1\). Then, in the following cases, the sequence of positions
\(\beta'_v=(p'_1,\dots,p'_m)\) is given as follows.

\begin{itemize}
\item[(1)] If
\[
\{p_1,\dots,p_m\}\cap\{k-1,k,k+1\}=\{k-1=p_t\},
\]
then
\[
p'_s=p_s \quad \text{for }s\neq t,
\qquad
p'_t=k.
\]

\item[(2)] If
\[
\{p_1,\dots,p_m\}\cap\{k-1,k,k+1\}
=
\{k-1=p_t,\ k=p_{t+1}\},
\]
then
\[
p'_s=p_s \quad \text{for }s\neq t,t+1,
\qquad
p'_t=k,\quad p'_{t+1}=k+1.
\]

\item[(3)] If
\[
\{p_1,\dots,p_m\}\cap\{k-1,k,k+1\}=\{k=p_t\},
\]
then
\[
p'_s=p_s \quad \text{for }s\neq t,
\qquad
p'_t=k-1.
\]

\item[(4)] If
\[
\{p_1,\dots,p_m\}\cap\{k-1,k,k+1\}
=
\{k=p_t,\ k+1=p_{t+1}\},
\]
then
\[
p'_s=p_s \quad \text{for }s\neq t,t+1,
\qquad
p'_t=k-1,\quad p'_{t+1}=k.
\]
\end{itemize}
\end{enumerate}
\end{lemma}

\begin{proof}
The first assertion is the standard left-descent criterion in Coxeter theory.
Indeed,
\[
s_i v<v
\]
if and only if \(s_i\) is a left descent of \(v\), which is equivalent to
\[
v^{-1}(\alpha_i)\in R^-.
\]

We prove the first two cases in (2). The last two follow by applying the same
argument to the inverse \(3\)-move.

\medskip
\noindent
\emph{First case.}
Assume
\[
\{p_1,\dots,p_m\}\cap\{k-1,k,k+1\}
=
\{k-1=p_t\}.
\]
Thus the local letter \(i_{k-1}=i\) is selected, while \(i_k=j\) and
\(i_{k+1}=i\) are not selected.

Set
\[
z:=v_{k-1}^{-1}v.
\]
Then
\[
v_{k-2}^{-1}v=s_i z.
\]
Since the letter at position \(k-1\) is selected, Lemma~\ref{lem:vk} gives
\[
z<s_i z.
\]
Equivalently,
\[
s_i z>z.
\]
Since the letter at position \(k\) is not selected, Lemma~\ref{lem:vk} gives
\[
s_j z>z.
\]
By part (1), we obtain
\[
z^{-1}(\alpha_i)\in R^+,
\qquad
z^{-1}(\alpha_j)\in R^+.
\]

In \(\beta'\), the local word is \((j,i,j)\). The candidate subexpression
obtained by replacing the selected position \(k-1\) in \(\beta\) by the
position \(k\) in \(\beta'\) still represents \(v\). Thus, by leftmostness,
\(\beta'_v\) is lexicographically no larger than this candidate.

We claim that \(\beta'_v\) cannot select position \(k-1\). Suppose otherwise.
Then, before position \(k-1\) in \(\beta'\), the remaining factor is \(s_i z\),
and selecting the local letter \(j\) gives, by Lemma~\ref{lem:vk},
\[
s_js_i z<s_i z.
\]
By part (1), this is equivalent to
\[
z^{-1}s_i(\alpha_j)\in R^-.
\]
Since \(c_{ij}=-1\), we have
\[
s_i(\alpha_j)=\alpha_i+\alpha_j.
\]
Therefore
\[
z^{-1}s_i(\alpha_j)
=
z^{-1}(\alpha_i)+z^{-1}(\alpha_j)\in R^-.
\]
This is impossible, because both \(z^{-1}(\alpha_i)\) and
\(z^{-1}(\alpha_j)\) are positive roots, and their sum is the root
\(z^{-1}s_i(\alpha_j)\). Hence \(p'_t\neq k-1\).

The candidate has \(p'_t=k\). Since \(\beta'_v\) is lexicographically no larger
than the candidate and cannot choose \(k-1\), it follows that
\[
p'_t=k.
\]

We next show that \(k+1\) is not selected. If \(p'_{t+1}=k+1\), then after the
position \(k\) has been selected, the remaining factor is \(z\), and selecting
the letter \(j\) at position \(k+1\) would imply
\[
s_jz<z
\]
by Lemma~\ref{lem:vk}. This contradicts \(s_jz>z\). Hence no further position
in \(\{k-1,k,k+1\}\) is selected.

Since \(\beta\) and \(\beta'\) coincide outside the local interval, the
remaining selected positions agree:
\[
p'_s=p_s\qquad(s\neq t).
\]
This proves the first case.

\medskip
\noindent
\emph{Second case.}
Assume
\[
\{p_1,\dots,p_m\}\cap\{k-1,k,k+1\}
=
\{k-1=p_t,\ k=p_{t+1}\}.
\]
Thus the local selected subword in \(\beta\) is \((i,j)\), and the position
\(k+1\) is not selected.

Set
\[
z:=v_k^{-1}v.
\]
Then
\[
v_{k-1}^{-1}v=s_jz,
\qquad
v_{k-2}^{-1}v=s_is_jz.
\]
Since the position \(k+1\) is not selected, Lemma~\ref{lem:vk} gives
\[
s_iz>z.
\]
By part (1), we have
\[
z^{-1}(\alpha_i)\in R^+.
\]

The candidate subexpression in \(\beta'\) selects the local positions
\[
k,\quad k+1,
\]
which again gives the local product \(s_is_j\). Hence the candidate has
\[
p'_t=k,\qquad p'_{t+1}=k+1.
\]

We claim that \(\beta'_v\) cannot select \(k-1\). Suppose otherwise. Then the
local letter at \(k-1\) in \(\beta'\) is \(j\), and before this position the
remaining factor is \(s_is_jz\). By Lemma~\ref{lem:vk}, selecting this letter
would imply
\[
s_js_is_jz<s_is_jz.
\]
By part (1), this gives
\[
z^{-1}s_js_i(\alpha_j)\in R^-.
\]
Since \(c_{ij}=-1\), we have
\[
s_js_i(\alpha_j)=\alpha_i.
\]
Thus
\[
z^{-1}(\alpha_i)\in R^-,
\]
contradicting \(z^{-1}(\alpha_i)\in R^+\). Hence \(p'_t\neq k-1\).

Since the candidate begins at \(k\), leftmostness forces
\[
p'_t=k.
\]
Moreover, the next selected position must be \(k+1\); otherwise the candidate
would be lexicographically smaller. Hence
\[
p'_{t+1}=k+1.
\]

Outside the local interval, \(\beta\) and \(\beta'\) coincide. Therefore
Lemma~\ref{lem:vk} gives
\[
p'_s=p_s
\qquad
\text{for all }s\neq t,t+1.
\]
This proves the second case.

The third and fourth cases are obtained by applying the first and second cases
to the inverse \(3\)-move
\[
(j,i,j)\longleftrightarrow(i,j,i).
\]
This completes the proof.
\end{proof}

\end{document}
